\documentclass{article}
\usepackage[utf8]{inputenc}
\usepackage{geometry}
\usepackage{graphicx}
\usepackage{pgfplots}
\usepackage{amsmath}
\usepackage{amssymb}
\usepackage{amsfonts}
\usepackage[thmmarks, amsmath, thref]{ntheorem}
\usepackage{mdframed}
\usepackage{amsthm}
\usepackage{physics}
\usepackage{proof}
\usepackage{derivative}
\usepackage{hyperref}
\usepackage{wrapfig}
\usepackage{amsthm}
\usepackage{xcolor}
\renewcommand{\theenumi}{\roman{enumi}}

\theoremstyle{nonumberplain}
\theoremheaderfont{\itshape}
\theorembodyfont{\upshape}
\theoremseparator{.}
\theoremsymbol{\ensuremath{\square}}
\theoremsymbol{\ensuremath{\blacksquare}}
\newtheorem{solution}{Solution}
\theoremseparator{. ---}
\theoremsymbol{\mbox{\texttt{;o)}}}

\pgfplotsset{compat=1.18}
\begin{document}

1. Si \( f:[a, b] \subset \mathbb{R} \rightarrow \mathbb{R} \) es clase \( C^{1} \); definimos la longitud de la gráfica de \( f \) en \( [a, b] \) como la longitud de la trayectoria \( t \mapsto(t, f(t)) \) con \( t \in[a, b] \). \\
(a) Pruebe que la longitud de la gráfica de \( f \) en \( |a, b| \) es
\[
\int_{a}^{b}\left(1+\left(\frac{\mathrm{d} f}{\mathrm{~d} x}\right)^{2}\right)^{\frac{1}{2}} \mathrm{~d} x
\]


Dividimos el intervalo \([a, b]\) en \( n \) subintervalos mediante una partición,
\[
a = x_0 < x_1 < \cdots < x_n = b
\]
Aproximamos la gráfica por la poligonal que une los puntos
\[
P_k = (x_k, f(x_k)), \quad k = 0, 1, \ldots, n
\]
La longitud de esta poligonal es:
\[
L_n = \sum_{k=1}^n \left\| P_k - P_{k-1} \right\| = \sum_{k=1}^n \sqrt{(x_k - x_{k-1})^2 + [f(x_k) - f(x_{k-1})]^2}
\]
Sacamos factor común \((x_k - x_{k-1})^2\):
\[
L_n = \sum_{k=1}^n (x_k - x_{k-1}) \sqrt{1 + \left( \frac{f(x_k) - f(x_{k-1})}{x_k - x_{k-1}} \right)^2 }
\]
Observamos que el cociente \(\frac{f(x_k) - f(x_{k-1})}{x_k - x_{k-1}}\) es el incremento medio de \(f\) en \([x_{k-1}, x_k]\), que por el Teorema del Valor Medio existe un punto \( \xi_k \in (x_{k-1}, x_k) \) tal que
\[
\frac{f(x_k) - f(x_{k-1})}{x_k - x_{k-1}} = f'(\xi_k)
\]
Así,
\[
L_n = \sum_{k=1}^n (x_k - x_{k-1}) \sqrt{1 + [f'(\xi_k)]^2 }
\]
Cuando el tamaño máximo de los subintervalos \( \max_k (x_k - x_{k-1}) \to 0 \), la suma anterior es una suma de Riemann para la función \( F(x) = \sqrt{1 + [f'(x)]^2} \) en \([a, b]\). Por la continuidad de \(f'\) (por ser de clase \(C^1\)), \(F\) es continua y la suma converge a la integral definida:
\[
L = \lim_{n \to \infty} L_n = \int_a^b \sqrt{1 + [f'(x)]^2} \, dx = \int_a^b \sqrt{1 + \left(\frac{df}{dx}\right)^2} \, dx
\]


(b) Encontrar la longitud de la gráfica de \( y=\ln (x) \) desde \( x=1 \) hasta \( x=2 \). \\

Observamos la raíz \(\sqrt{x^2 + 1}\). Esto sugiere una sustitución trigonométrica:
\[
x = \tan\theta, \quad \theta \in \left[ \arctan 1, \arctan 2 \right]
\]
Entonces:
\[
dx = \sec^2\theta\, d\theta
\]
\[
\sqrt{x^2 + 1} = \sqrt{\tan^2\theta + 1} = \sqrt{\sec^2\theta} = \sec\theta
\]
Los límites cambian como sigue...\\

Cuando \(x=1\), \(\theta_1 = \arctan 1 = \frac{\\pi}{4}\)\\
Cuando \(x=2\), \(\theta_2 = \arctan 2\)\\

Reescribimos la integral:
\[
I = \int_{\theta_1}^{\theta_2} \frac{\sec\theta}{\tan\theta} \cdot \sec^2\theta\, d\theta
\]
Recordemos que:
\[
\tan\theta = \frac{\sin\theta}{\cos\theta}, \quad \sec\theta = \frac{1}{\cos\theta}
\]
Por tanto:
\[
\frac{\sec\theta}{\tan\theta} = \frac{1/\cos\theta}{\sin\theta/\cos\theta} = \frac{1}{\sin\theta}
\]
Así que:
\[
I = \int_{\theta_1}^{\theta_2} \frac{1}{\sin\theta} \cdot \sec^2\theta\, d\theta = \int_{\theta_1}^{\theta_2} \frac{1}{\sin\theta} \cdot \frac{1}{\cos^2\theta} d\theta = \int_{\theta_1}^{\theta_2} \frac{1}{\sin\theta\cos^2\theta} d\theta
\]

Obsérvese que:
\[
\frac{1}{\sin\theta \cos^2\theta} = \frac{1}{\cos^2\theta} \cdot \frac{1}{\sin\theta} = \sec^2\theta \csc\theta
\]
Vamos a integrar \(\sec^2\theta \csc\theta\):

Sea \(u = \csc\theta\), \(dv = \sec^2\theta d\theta\):
\[
u = \csc\theta \implies du = -\csc\theta \cot\theta d\theta
\]
\[
dv = \sec^2\theta d\theta \implies v = \tan\theta
\]
Por integración por partes:
\[
\int u\, dv = uv - \int v\, du
\]
\[
I = \csc\theta \tan\theta - \int \tan\theta \left(-\csc\theta \cot\theta\right) d\theta
\]
\[
= \csc\theta \tan\theta + \int \tan\theta \csc\theta \cot\theta d\theta
\]
Ahora, simplificamos \(\tan\theta \csc\theta \cot\theta\):
\[
\tan\theta = \frac{\sin\theta}{\cos\theta}
\]
\[
\csc\theta = \frac{1}{\sin\theta}
\]
\[
\cot\theta = \frac{\cos\theta}{\sin\theta}
\]
Por tanto:
\[
\tan\theta \csc\theta \cot\theta = \frac{\sin\theta}{\cos\theta} \cdot \frac{1}{\sin\theta} \cdot \frac{\cos\theta}{\sin\theta} = \frac{1}{\cos\theta} \cdot \frac{\cos\theta}{\sin\theta} = \frac{1}{\sin\theta}
\]
Así que:
\[
I = \csc\theta \tan\theta + \int \frac{1}{\sin\theta} d\theta
\]
\[
= \csc\theta \tan\theta + \int \csc\theta d\theta
\]
Recordemos que:
\[
\int \csc\theta d\theta = \ln \left| \tan\left(\frac{\theta}{2}\right) \right| + C
\]

4. Reescribiendo la integral
Por lo tanto:
\[
I = \csc\theta \tan\theta + \ln \left| \tan\left(\frac{\theta}{2}\right) \right| + C
\]
Ahora, necesitamos volver a la variable original \(x\):

\(x = \tan\theta \implies \theta = \arctan x\)\\
\(\tan\theta = x\)\\
\(\sin\theta = \frac{x}{\sqrt{x^2 + 1}}\)\\
\(\cos\theta = \frac{1}{\sqrt{x^2 + 1}}\)\\
\(\csc\theta = \frac{1}{\sin\theta} = \frac{\sqrt{x^2 + 1}}{x}\)\\

Entonces:
\[
\csc\theta \tan\theta = \left(\frac{\sqrt{x^2+1}}{x}\right) \cdot x = \sqrt{x^2+1}
\]
Queda por expresar \(\tan\left( \frac{\theta}{2} \right)\) en términos de \(x\), para ello sabemos que si \(x = \tan\theta\), entonces
\[
\tan\left(\frac{\theta}{2}\right) = \frac{x}{1 + \sqrt{x^2+1}}
\]
Esta es una identidad trigonométrica estándar, por lo tanto:
\[
I = \sqrt{x^2+1} + \ln \left| \frac{x}{1+\sqrt{x^2+1}} \right| + C
\]

Recordemos que
\[
I = \int_{1}^{2} \frac{\sqrt{x^2+1}}{x} dx = \left[ \sqrt{x^2+1} + \ln \left( \frac{x}{1+\sqrt{x^2+1}} \right) \right]_{x=1}^{x=2}
\]

Así que el resultado es:
\[
I = \left[ \sqrt{5} + \ln \left( \frac{2}{1+\sqrt{5}} \right) \right] - \left[ \sqrt{2} + \ln \left( \frac{1}{1+\sqrt{2}} \right) \right] \approx 1.22201
\]

\end{document}



